\begin{figure}[t]
  \centering
  \begin{subfigure}{0.24\textwidth}
    \includegraphics[width=\linewidth]{ch07/equivariance-symmetric-m8-t0.pdf}
    \caption{$t=0$}
  \end{subfigure}\hfill
  \begin{subfigure}{0.24\textwidth}
    \includegraphics[width=\linewidth]{ch07/equivariance-symmetric-m8-t1.pdf}
    \caption{$t=T/3$}
  \end{subfigure}\hfill
  \begin{subfigure}{0.24\textwidth}
    \includegraphics[width=\linewidth]{ch07/equivariance-symmetric-m8-t2.pdf}
    \caption{$t=2T/3$}
  \end{subfigure}\hfill
  \begin{subfigure}{0.24\textwidth}
    \includegraphics[width=\linewidth]{ch07/equivariance-symmetric-m8-t3.pdf}
    \caption{$t=T$}
  \end{subfigure}
  \caption{Permutation-symmetric multi-agent configuration with $M=8$ interchangeable drones. Antipodal initial and terminal positions on a circle provide a canonical setting for permutation equivariance; predicted collision-free trajectories at four time snapshots are shown in \cref{sec:oc:2d-experiments}.}
  \label{fig:oc:equivariance}
\end{figure}
